\documentclass[10pt,a4paper]{amsart}
\usepackage[margin=19mm]{geometry}
\usepackage{amsmath,amssymb,mathtools}
\usepackage[scr=rsfs,cal=euler]{mathalfa}
\usepackage{xfrac}
\usepackage[unicode,hidelinks,bookmarks=false]{hyperref}
\newcommand{\ie}{i.e.\ }
\newcommand{\cf}{cf.\ }
\newcommand{\resp}{respectively\ }
\newcommand{\Subsection}{\subsection}
\providecommand{\nobreakdash}{\nobreak\hbox{-}}
\setlength{\emergencystretch}{2em}
\multlinegap0pt
\usepackage{bm}
\usepackage{bbm}
\usepackage{dsfont}
\usepackage{xspace}
\usepackage{cancel}
\usepackage{mathtools}
\usepackage{amsthm}
\makeatletter
\@ifundefined{proposition}{\newtheorem{proposition}{Proposition}}{}
\makeatother
\newcommand{\al}{\alpha}
\newcommand{\be}{\beta}
\newcommand{\dpr}[2]{\langle #1,#2 \rangle}
\newcommand{\eps}{\epsilon}
\newcommand{\f}[2]{\frac{#1}{#2}}
\newcommand{\ga}{\gamma}
\newcommand{\la}{\lambda}
\newcommand{\p}{\partial}
\newcommand{\rone}{\mathbf R}
\newcommand{\si}{\sigma}
\title[Asymptotic stability for monotone traveling kinks of the dis]{Asymptotic stability for monotone traveling kinks of the dissipative Boussinesq problem}
\author{Atanas G. Stefanov}
\date{Source: arXiv:2606.27690v1 (CC BY 4.0)}
\begin{document}
\maketitle

\noindent\emph{Source and licence.} Asymptotic stability for monotone traveling kinks of the dissipative Boussinesq problem. Version arXiv:2606.27690v1 (CC BY 4.0), distributed under the Creative Commons Attribution 4.0 International licence, as recorded in the arXiv licence metadata for this exact version and shown on the abstract page https://arxiv.org/abs/2606.27690v1. Full rights notice: licenses/arxiv-2606.27690-CC-BY-4.0.txt.\par\smallskip

\noindent\textit{Editorial scope.} Counted unit: the Proposition whose source label is \texttt{prop:65} in the article (source0.tex lines 691--705, source label \texttt{prop:65}), delivered together with the article's own complete original proof of it (source0.tex lines 706--854, one \texttt{proof} environment of 149 lines). No mathematics is written, repaired or re-derived: statement and proof are the article's own text line for line. Required context retained uncounted: f:1 (lines 409--416); 105 (lines 501--508); 400 (lines 564--567); 420 (lines 603--608); prop:54 (lines 639--650). Every definition, display and label of those lines is retained unchanged. Omitted from this package and not redistributed: 1--408, 417--500, 509--563, 568--602, 609--638, 651--690, 855--1086. Nothing inside a retained excerpt is omitted or altered: every retained body is delivered exactly as the article's lines are written, so no omission receipt of any kind accompanies this package. Citations: the retained excerpts carry no citation command at all, so no bibliography entry of the article is transcribed and no reference is redistributed. Reference adaptation: none was needed and none was made; every reference inside the delivered unit resolves inside it, so no unresolved reference or citation is produced. Printed slips kept literal: no printed word, symbol, hypothesis or number of the counted statement or of its proof was altered. Layout and class: the article's own class is \texttt{amsart} with packages including geometry, amsmath, amsfonts, amsthm, amssymb, bbm, graphicx, color, dsfont, fourier; the wrapper is the lane's pinned \texttt{amsart} editorial wrapper at 10pt on A4, which restates the theorem-like environments the retained text needs and adds the article's own macro definitions that the retained text needs (\texttt{\textbackslash al}, \texttt{\textbackslash be}, \texttt{\textbackslash dpr}, \texttt{\textbackslash eps}, \texttt{\textbackslash f}, \texttt{\textbackslash ga}, \texttt{\textbackslash la}, \texttt{\textbackslash p}, \texttt{\textbackslash rone}, \texttt{\textbackslash si}), with extra wrapper packages (bm, bbm, dsfont, xspace, cancel, mathtools). The delivered numbering is the unit's own. Assets: no retained span contains a figure environment, a graphics-inclusion command, a PDF-page inclusion, a table or a TikZ picture; no image file is copied into this package, the package declares no asset dependency and no image file is redistributed with it. Licence: version arXiv:2606.27690v1 of this article is distributed under the Creative Commons Attribution 4.0 International licence, as recorded in the arXiv licence metadata for this exact version and shown on the abstract page https://arxiv.org/abs/2606.27690v1. A full rights notice is delivered at licenses/arxiv-2606.27690-CC-BY-4.0.txt. Characters: every retained excerpt is pure ASCII, so the delivered engine, XeLaTeX with its default Latin Modern text font, reports no missing character.
\begin{eqnarray}
	\label{f:1} 
	& & \begin{cases}
		\phi(x)=\phi_\pm - c_\pm exp(\f{m \mu_{\pm}(\nu)}{\si} x)+o(e^{\f{m \mu_{\pm}(\nu)}{\si} x}), \ \ x\sim \pm \infty \\
		\phi'(x)=- c_\pm \f{m \mu_{\pm}(\nu)}{\si} exp(\f{m \mu_{\pm}(\nu)}{\si} x)+o(e^{\f{m \mu_{\pm}(\nu)}{\si} x}), \ \ x\sim \pm \infty \\
		\phi''(x)=- c_\pm \left(\f{m \mu_{\pm}(\nu)}{\si}\right)^2 exp(\f{m \mu_{\pm}(\nu)}{\si} x)+o(e^{\f{m \mu_{\pm}(\nu)}{\si} x}), \ \ x\sim \pm \infty 
		\end{cases}. 
\end{eqnarray}

\begin{equation}
	\label{105} 
	\begin{cases}
		U_t +(\nu'(t)-\si)  U'+\mu' \phi' - \al U''= V' \\
		V_t+(\nu'(t)-\si) V' +\mu'\psi'-\be  V''  = -\ga  U''' + U' + \f{1}{2} (2\phi U+U^2)', \\
		U(0,y)=U_0(y) \in L^2(\rone) ; V(0,y)= V_0(y) \in L^2(\rone). 
	\end{cases}
	\end{equation} 

\begin{equation}
	\label{400} 
	w_{tt} - w_{txx}+ B w_{xxxx} - w_{xx}= \f{1}{2} \p_{xx} (2\phi(\cdot-\si t)w+w^2).
\end{equation}

\begin{eqnarray}
	\label{420} 
		\hat{q}(t, \xi) &=& \f{\la_-(\xi) e^{t \la_+(\xi)}-\la_+(\xi) e^{t \la_-(\xi)}}{\la_-(\xi)-\la_+(\xi)} \hat{q}(0,\xi) + \f{e^{t \la_+(\xi)}-e^{t \la_-(\xi)}}{\la_+(\xi)-\la_-(\xi)} \hat{q}_t(0,\xi) \\
		\nonumber
		& & -4\pi^2 \int_0^t \hat{G}(t-s, \xi) [\xi^2 \hat{N}(s, \xi)] ds,
\end{eqnarray}

\begin{proposition}
	\label{prop:54} 
	Suppose that $w_0\in H^1(\rone), w_1\in L^2(\rone)$. Then, there exists time \\  $T=T((\|w_0\|_{H^1(\rone)}, \|w_1\|_{L^2(\rone)})>0$ and a unique solution 
	$w\in C([0,T), H^1(\rone))$ of \eqref{400}. More specifically, $w$ satisfies the integral equation \eqref{420}, with the nonlinearity $N(t,x)=2\phi(x-\si t)w(t,x) +w^2(t,x)$. 
	A posteriori, the solution $w$ described above is in fact infinitely smooth, $w\in C^\infty([0,T), H^\infty(\rone))$. 
	
	The blow-up alternative for \eqref{400} is as follows - the solution $w$ blows up at time $T^*$ if and only if 
	\begin{equation}
		\label{90P}
		\limsup_{t\to T^*-} \|w(t, \cdot)\|_{H^1(\rone)}=+\infty. 
	\end{equation}
\end{proposition}

\begin{proposition}
	\label{prop:65} 
There exists $\eps_0=\eps_0(\phi, \al, \be, \ga)>0$, so that whenever 
\begin{equation}
	\label{297} 
	U_0\in H^1(\rone), V_0\in L^2(\rone):  \|U_0\|_{H^1}+\|V_0\|<\eps_0, 
\end{equation}
the solution $U,V$ to \eqref{105} extends globally in time. Moreover, they are smooth functions, and they satisfy the estimates 
\begin{eqnarray}
	\label{w:10} 
& & 	\|U(t)\|_{H^1}+ \|V(t)\|\leq C_1(\phi) (\|U_0\|_{H^1}+\|V_0\|) \\
		\label{w:20}
	& & 	\int_0^\infty  [\|U'(t)\|^2+ \|U''(t)\|^2+\|V'(t)\|^2]dt\leq C_2(\phi) (\|U_0\|_{H^1}^2+\|V_0\|^2). 
\end{eqnarray}
\end{proposition}

\begin{proof} 
	First, a word about the strategy of the proof. We run the usual continuity argument. Specifically,  we start the small data as in \eqref{297}, then analyze the local solutions, which exist thanks to  Proposition \ref{prop:54}. We need to show that \eqref{w:10} persists forever, where $C_1(\phi)$ will be determined explicitly in the course of the argument, see \eqref{419} below.  The trick is to use \\  $\sup_{0<t<t_*}	\|U(t)\|_{H^1}+ \|V(t)\|\leq C_1(\phi)  (\|U_0\|_{H^1}+\|V_0\|)$ as an {\it a priori} bound and show that until it is satisfied, which maybe up to some finite time  $t_*$, we have in fact a better bound, say $\limsup_{t\to t_*-} \|U(t)\|_{H^1}+ \|V(t)\|\leq \f{C_1(\phi)}{2} (\|U_0\|_{H^1}+\|V_0\|)$. This argument shows that the bound $\sup_{0<t<t_*}	\|U(t)\|_{H^1}+ \|V(t)\|\leq C_1(\phi)  (\|U_0\|_{H^1}+\|V_0\|)$ persists for all $t_*<\infty$. 
	
Let us now proceed to the actual estimates. Consider  \eqref{105} and  take dot product of the first equation with $U$ and the second equation with $V$, and add them up. Recall however that $\phi=\phi(y+\mu(t))$, whereas $U=U(y+\nu(t),t), V=V(y+\nu(t),t)$.
We obtain 
\begin{equation}
\label{110} 
\begin{cases}
	\f{1}{2} \p_t\left( \|U\|^2+\|V\|^2\right)+\mu'(\dpr{\phi'}{U}+\dpr{\psi'}{V})+ \al \|U'\|^2+\be\|V'\|^2= \\
	= -\ga\dpr{U'''}{V}+\f{1}{2} \dpr{\p_x (2\phi U+U^2)}{V}.
\end{cases}	
\end{equation}
We now compute various terms, by taking into account the relation \eqref{105}, 
\begin{eqnarray*}
	-\ga\dpr{U'''}{V} &=& \ga \dpr{U''}{V'}= \ga \dpr{U''}{U_t+(\nu'-\si)   U'+\mu'\phi'-\al U''}= \\
	&=& -\f{\ga}{2} \p_t \|U'\|^2+\ga\mu'\dpr{U}{\phi'''}-\ga \al \|U''\|^2. 
\end{eqnarray*}
On the other hand, integration by parts yields 
\begin{eqnarray*}
& & 	\f{1}{2}  \dpr{\p_x (2\phi U+U^2)}{V} = -\dpr{\phi U}{V'}-\f{1}{2}  \dpr{U^2}{V'}=  \\
	&=& - \dpr{\phi U}{U_t+(\nu'-\si) U'+\mu'\phi'-\al U''}- \f{1}{2}  \dpr{U^2}{U_t- (\nu'-\si)   U'+\mu'\phi'-\al U''}. 
\end{eqnarray*}
Term by term, and integration by parts 
\begin{eqnarray*}
	- \dpr{\phi U}{U_t} & = & -\f{1}{2} \p_t \left( \int \phi U^2 \right) +\f{1}{2}  \mu' \int \phi' U^2. \\
	(\si-\nu')  \dpr{\phi U}{U'} &=& \f{\nu'-\si}{2} \int \phi' U^2; \\
	- \mu' \dpr{\phi U}{\phi'}  &=&    -\mu' \int \phi\phi' U, \\
\al  \dpr{\phi U}{U''} &=& - \al \int \phi (U')^2 + \f{\al}{2}  \int \phi'' U^2. \\
-\f{1}{2} \dpr{U^2}{U_t} &=& -\f{1}{6} \p_t \int U^3(t,y) dy,\\
-\f{1}{2} \mu'\dpr{U^2}{\phi'} &=& -\f{1}{2} \mu'\int \phi' U^2,       \\ 
\f{\al}{2}   \dpr{U^2}{U''} &=& - \al \int  U (U')^2, 
\end{eqnarray*}
	All in all, we arrive at the identity, 
	\begin{eqnarray*}
	& &	\f{1}{2} \p_t\left( \int (1+ \phi) U^2+\f{1}{3} \int U^3+\|V\|^2+ \ga \|U'\|^2 \right)+\\
	& & + \al (\|U'\|^2+\int \phi (U')^2) + \be\|V'\|^2 + \ga \al \|U''\|^2+ \al \int  U (U')^2+ \\
& & + \mu'\left(\dpr{\phi'-\ga \phi'''+  \phi\phi'}{U} +\dpr{\psi'}{V}\right)+\f{1}{2}  \int  ((\si-\nu')  \phi' - \al \phi'')  U^2  =0.
	\end{eqnarray*}
 Here, the important terms to take care first are on the last line. More precisely, we need to make an appropriate selection of the functions $\mu, \nu$. First, we shall select $\nu(t):=N t+\si t$ for sufficiently large $N=N(\phi)$, to be determined in a moment. 
 
 The next step is to select $\mu'$, so that $\mu'\left(\dpr{\phi'-\ga \phi'''+  \phi\phi'}{U} +\dpr{\psi'}{V}\right)\ge 0$. Thus, the natural choice  is to take  
 \begin{equation}
 	\label{300}
 	\mu'(t)= \dpr{\phi'-\ga \phi'''+  \phi\phi'}{U} +\dpr{\psi'}{V}, \ \ \mu(0)=0. 
 \end{equation}
% \begin{equation}
% 	\label{300} 
% 	\begin{cases}
% 		
% 		=\int_{-\infty}^{+\infty} [(\phi'-\ga \phi'''+\phi\phi')(y+\mu(t)-\nu(t)) U(t,y) dy +    \psi'(y+\mu(t)-\nu(t)) V(t, y)] dy.
% 	\end{cases}
% \end{equation}
 Note that, 
 	\begin{eqnarray*}
 \dpr{\phi'-\ga \phi'''+  \phi\phi'}{U} +\dpr{\psi'}{V} &=&  \int_{-\infty}^{+\infty} (\phi'-\ga \phi'''+\phi\phi')(y+\mu(t)-\nu(t)) U(t,y) dy +  \\
 &+& 
\int_{-\infty}^{+\infty}  \psi'(y+\mu(t)-\nu(t)) V(t, y) dy=:G(t)
	\end{eqnarray*}
So, {\it assuming that we have control over $\|U(t, \cdot)\|, \|V(t, \cdot)\|$} till some time, say $t_*$, then the ODE for $\mu$ does not blow-up. In fact 
\begin{equation}
	\label{3001}
	\sup_{0<t<t_*}|\mu'(t)|=\sup_{0<t<t_*} |G(t)|\leq C \sup_{0<t<t_*}(\|U(t, \cdot)\|+ \|V(t, \cdot)\|),
\end{equation}
where $C=C(\phi)$. 

% 
%Indeed, \eqref{300} is an ODE for $\mu$, in the form $\mu'(t)=G(t, \mu(t))$, where $G$ is Lipschitz in the $\mu$ variable, due to the smoothness of $\phi$. Also, in order to justify the construction of $\mu(t)$, we shall need to establish that $G$ is continuous in the $t$ variable, which would follow from the Lipschitzness of $G$. To this end, it will suffice to establish that the following two functions 
%\begin{equation}
%	\label{310} 
%	t\to \int_{-\infty}^{+\infty} \psi'(y+\mu-\nu(t)) V_t(t, y) dy, t\to \int_{-\infty}^{+\infty} [(\phi'-\ga \phi'''+\phi\phi')(y+\mu-\nu(t)) U_t(t,y) dy
%\end{equation}
%are uniformly bounded up to a potential blow-up time $t_*$.   {\it That is, \eqref{310} still needs to be established in a pending continuity in time argument, see below}. 
With the choice of $\nu(t), \mu(t)$ as made in \eqref{300},  we arrive at the energy estimate 
\begin{equation}
	\label{320} 
	\begin{cases}
		\f{1}{2}  \p_t\left( \int (1+ \phi) U^2+\f{1}{3} \int U^3+\|V\|^2+ \ga \|U'\|^2 \right) + \al \int  U (U')^2\\
		+  \be\|V'\|^2 +  \ga \al \|U''\|^2+ \al  \int (1+ \phi) (U')^2 +\f{1}{2}  \int  (N (-\phi') - \al \phi'')  U^2  \leq 0.
	\end{cases}
\end{equation} 
Here, we have the easy Sobolev bound for the cubic term (for some absolute constant $C$), 
$$
 \left|\int  U (U')^2\right|  \leq \|U\|_{L^\infty} \|U'\|_{L^2}^2\leq C \|U(t)\|_{H^1} \|U'\|_{L^2}^2\leq C  C_1(\phi) \eps_0\|U'\|_{L^2}^2,
$$
due to the {\it a priori} bound on $\sup_{0<t<t_*} \|U\|_{H^1}\leq C_1(\phi)\eps_0$ and \eqref{297}. 
Additionally, we have the bound 
$$
 \f{\al}{2}   \int (1+ \phi) (U')^2 +\al  \int  U (U')^2\geq  \left(\f{\al(1+\phi_+)}{2} -CC_1(\phi) \eps_0\right)  \int (U')^2 >0, 
$$
for small enough $\eps_0$, namely $CC_1(\phi) \eps_0<\f{\al(1+\phi_+)}{4} $. 

We now tackle the remaining quadratic in $U$ terms, which appears in \eqref{320}. 
%$$
% \f{\al}{4}   \int (1+ \phi) (U')^2 +\f{1}{2}  \int  (N (-\phi') - \al \phi'')  U^2\geq  \f{1}{2}  \int  (N (-\phi') - \al \phi'')  U^2. 
%$$
 Recalling that the kink $\phi$ is decreasing in $(-\infty, +\infty)$, and hence $-\phi'>0$, we may select $N=N(\phi)$, so that $N (-\phi') - \al \phi''=(-\phi')\left(N-\al\f{\phi''}{-\phi'}\right)>0$ and hence $$
 \int  (N (-\phi') - \al \phi'')  U^2>0.
 $$
  Indeed, the asymptotics \eqref{f:1} show that 
$$
\lim_{x\to \pm \infty} \left|\frac{\phi''(x)}{\phi'(x)}\right|= \f{m|\mu_{\pm}(\nu)|}{\si}>0, 
$$
whence it follows,  that $\sup_{x} \left|\frac{\phi''(x)}{\phi'(x)}\right|<\infty$, as $\phi'$ does not vanish. So, a choice of $N:=2\al \sup_{x} \left|\frac{\phi''(x)}{\phi'(x)}\right|$ ensures that $N (-\phi') - \al \phi''>0$. 
 


It follows that the last quadratic form is non-negative, which means that 
\begin{equation}
	\label{330} 
	\f{1}{2} \p_t\left( \int (1+ \phi) U^2+\f{1}{3} \int U^3+\|V\|^2+ \ga \|U'\|^2 \right) +  \be\|V'\|^2 +\f{\al(1+\phi_+)}{2} \|U'\|^2+  \ga \al \|U''\|^2\leq 0.
\end{equation}
This means that up to  $t_*$, we have control of the various norms, per \eqref{330} as follows. Ignoring the positive terms $ \be\|V'\|^2 +\f{\al(1+\phi_+)}{2} \|U'\|^2+  \ga \al \|U''\|^2$ for the moment, we  integrate\eqref{330}, to obtain for all $0<t<t_*$, 
\begin{eqnarray*}
& & 	 \int  (1+\phi(y+\mu(t)-\nu(t)) ) U^2(t, y)+\f{1}{3} \int U^3(t,y)+\|V(t,\cdot)\|^2+ \ga \|U'(t,\cdot)\|^2 \leq \\
	 &\leq &  \int  (1+\phi(y))  U_0^2(t, y)+\f{1}{3} \int U_0^3(y)+\|V_0\|^2+ \ga \|U'_0\|^2. 
\end{eqnarray*}
Trivially, for an absolute constant $C$, 
$$
|\int_{-\infty}^{+\infty}  U^3(t,y) dy|\leq C \|U\|_{L^\infty} \|U\|^2\leq C \|U(t)\|_{H^1}^3\leq C_1(\phi) \eps_0 \|U(t)\|_{H^1}^2, 
$$
due to the {\it a priori} bound for $\|U(t)\|_{H^1}$. 
For $\eps_0: C_1(\phi)\eps_0<\f{1}{100} \min(\ga, 1+\phi_+)$, we can hide the contribution of $\int U^3$. Indeed, 
\begin{eqnarray*}
	& & \f{\ga}{2} \|U'(t,\cdot)\|^2+  \f{1}{2} \int  (1+\phi(y+\mu(t)-\nu(t)) ) U^2(t, y) + \f{1}{3} \int U^3(t,y)\geq \\
	&\geq & \left(\f{\ga}{2}+\f{1+\phi_+}{2}\right) \|U\|_{H^1}^2 - C_1(\phi) \eps_0 \|U(t)\|_{H^1}^2\geq 0, 
\end{eqnarray*}
by the choice of $\eps_0$. So, we arrive at, 
\begin{equation}
	\label{414} 
	\f{\min((1+\phi_+), \ga)}{2} \|U(t, \cdot)\|_{H^1}^2 + \|V(t,\cdot)\|^2 \leq 	\max((1+\phi_-),\ga)  \|U_0\|_{H^1}^2+\|V_0\|^2. 
\end{equation}
It follows that 
\begin{equation}
	\label{418} 
	\sup_{0<t<t_*} (\|U(t, \cdot)\|_{H^1}+\|V(t,\cdot)\|)< \f{C_1(\phi)}{2}( \|U_0\|_{H^1}+\|V_0\|),
\end{equation}
where 
\begin{equation}
	\label{419} 
C_1(\phi)=100\sqrt{ \f{1+\max((1+\phi_-),\ga)}{\min(\f{\min((1+\phi_+), \ga)}{2},1)}}. 
\end{equation}
  Going back to \eqref{330}, we see that we get additional information. Namely, adding back on the terms $\be\|V'\|^2 +\f{\al(1+\phi_+)}{2} \|U'\|^2+  \ga \al \|U''\|^2$ and integrating (having in mind that the rest of the terms are positive and may be ignored), we get an $L^2_t$ estimate 
\begin{equation}
	\label{360} 
	\int_0^{\infty} [\|U'(t, \cdot)\|^2+\|V'(t, \cdot)\|^2 + \|U''(t,\cdot)\|^2] dt\leq C  [\|U_0\|_{H^1}^2 +
	\|V_0\|_{L^2}^2].
\end{equation}
This completes the proof of Proposition \ref{prop:65}. 
\end{proof}

\end{document}
